\documentclass[11pt,a4paper]{article}

\usepackage[utf8]{inputenc}
\usepackage[T1]{fontenc}
\usepackage[spanish,es-nodecimaldot]{babel}
\usepackage{geometry}
\usepackage{array}
\usepackage{longtable}
\usepackage{enumitem}
\usepackage{xcolor}
\usepackage{hyperref}

\geometry{margin=2.4cm}
\hypersetup{
  colorlinks=true,
  linkcolor=blue!55!black,
  urlcolor=blue!55!black
}

\setlength{\parskip}{0.55em}
\setlength{\parindent}{0pt}
\setlist[itemize]{topsep=0.25em,itemsep=0.18em,leftmargin=1.4em}
\setlist[enumerate]{topsep=0.25em,itemsep=0.18em,leftmargin=1.6em}

\newcommand{\sesion}[1]{\textbf{Sesión #1}}
\newcommand{\producto}{\textbf{Producto verificable:} }
\newcommand{\aidriven}{\textit{AI-driven}}

\title{\textbf{Plan de Taller de 4 Semanas}\\
\large Postproceso RF, DSP aplicado y cumplimiento normativo paso a paso}
\author{Propuesta de formación técnica}
\date{2026-05-21}

\begin{document}

\maketitle

\section*{Resumen ejecutivo}

Este documento propone un taller de 4 semanas para que un equipo técnico con poca experiencia en DSP aprenda a comprender, diseñar, probar y operar el postproceso de una plataforma de sensado espectral. El foco no está en leer una base de código existente línea por línea, sino en dominar los conceptos que hacen que ese software tenga sentido: adquisición IQ, estimación espectral, ruido, umbrales, detección de emisiones, medición de ancho de banda, corrección de artefactos, contratos de datos, servicios de análisis y cumplimiento normativo.

El taller se apoya en dos dominios:

\begin{itemize}
  \item \textbf{Procesamiento en el nodo embebido:} captura de muestras IQ, almacenamiento temporal de datos, cálculo de PSD, compensación básica de errores, filtrado, demodulación opcional y preparación de tramas espectrales.
  \item \textbf{Procesamiento en plataforma:} recepción de una trama espectral, detección de emisiones, cálculo de parámetros, comparación contra información normativa, generación de resultados y entrega de respuesta a otros módulos de la plataforma.
\end{itemize}

El enfoque será \aidriven: se usará inteligencia artificial como apoyo para explicar conceptos, generar ejemplos de señales, crear casos de prueba, revisar unidades, documentar decisiones y preparar guías de diagnóstico. La IA será una herramienta de aprendizaje y productividad, no una autoridad técnica que sustituya la validación humana.

Al terminar, los participantes deberán poder explicar el flujo completo:

\begin{quote}
señal RF $\rightarrow$ muestras IQ $\rightarrow$ espectro/PSD $\rightarrow$ trama de medición $\rightarrow$ detección de emisiones $\rightarrow$ análisis normativo $\rightarrow$ resultado para la plataforma.
\end{quote}

\section*{Propósito del taller}

El propósito es llevar a los participantes desde una comprensión inicial de DSP hasta una capacidad práctica para razonar el postproceso completo. Se busca que puedan operar y extender el sistema con criterio: saber qué significa cada dato, qué supuestos existen, qué pruebas conviene hacer y qué riesgos aparecen al interpretar potencia, ruido, frecuencia y cumplimiento.

Al finalizar, el equipo debe poder:

\begin{itemize}
  \item Explicar IQ, FFT, PSD, RBW, ventana, ruido, umbral, DC, PPM, OBW, SNR y dBm en un lenguaje operativo.
  \item Entender cómo un nodo de adquisición convierte señales RF en un espectro utilizable.
  \item Diseñar el contrato de datos entre nodo de adquisición y servicio de análisis.
  \item Diferenciar potencia relativa digital, potencia integrada y estimaciones de intensidad de campo.
  \item Probar detección de emisiones con señales sintéticas.
  \item Ajustar y explicar umbrales de detección sin tratarlos como números mágicos.
  \item Entender el flujo de cumplimiento normativo: frecuencia, ancho de banda, potencia, licencia y tolerancias.
  \item Diseñar pruebas de borde para ruido, emisiones débiles, señales cercanas, errores de unidades y datos incompletos.
  \item Preparar una guía de operación y diagnóstico para el postproceso.
\end{itemize}

\section*{Audiencia objetivo}

El taller está dirigido a perfiles técnicos que tienen bases de programación o integración, pero no necesariamente formación profunda en procesamiento digital de señales:

\begin{itemize}
  \item Ingenieros de sistemas, electrónicos, telecomunicaciones o perfiles equivalentes.
  \item Tecnólogos o desarrolladores con bases de Python, terminal y JSON.
  \item Integradores que conectan sensores, backend, base de datos y servicios externos.
  \item Personal técnico encargado de operar, validar o interpretar reportes de cumplimiento.
\end{itemize}

No se asume dominio previo de DSP. Sí se recomienda:

\begin{itemize}
  \item Programación básica.
  \item Uso de terminal y Git.
  \item Nociones de HTTP, JSON y servicios web.
  \item Capacidad para leer gráficas y tablas.
  \item Interés por radiofrecuencia y medición espectral.
\end{itemize}

\section*{Enfoque pedagógico}

El taller se organiza como una progresión desde intuición hacia operación:

\begin{itemize}
  \item Primero se trabaja con señales pequeñas y visuales.
  \item Luego se introducen conceptos DSP sólo al nivel necesario.
  \item Después se diseña el flujo de datos como contrato de software.
  \item Finalmente se conecta detección, medición y cumplimiento.
\end{itemize}

La profundidad matemática será controlada. Se explicarán las ideas necesarias para tomar decisiones técnicas, pero no se convertirá el taller en una clase extensa de teoría de señales.

\begin{itemize}
  \item Se usará simulación antes de depender de hardware.
  \item Se trabajará con diagramas, gráficas, tablas de unidades y casos de prueba.
  \item Se separará el procesamiento del nodo de adquisición del procesamiento de plataforma.
  \item Se insistirá en trazabilidad: de dónde viene cada dato, en qué unidad está y para qué se usa.
  \item Se evitarán referencias a implementación puntual, detalles internos o rutas de código.
  \item Se enseñará a validar resultados con pruebas, no sólo a confiar en salidas automáticas.
\end{itemize}

\section*{Rol de la inteligencia artificial}

El enfoque \aidriven{} busca que los participantes usen IA para acelerar aprendizaje, documentación y pruebas. La IA debe ayudar a convertir conceptos técnicos en ejemplos ejecutables, listas de verificación y explicaciones comparables.

\subsection*{Usos esperados de IA}

\begin{itemize}
  \item Explicar conceptos DSP con analogías y ejemplos numéricos simples.
  \item Generar señales sintéticas para práctica.
  \item Proponer casos de prueba para detección y cumplimiento.
  \item Revisar consistencia de unidades en tablas y contratos de datos.
  \item Ayudar a redactar diagramas de flujo y guías de operación.
  \item Preparar checklists de diagnóstico.
  \item Sugerir preguntas para validar supuestos técnicos.
  \item Ayudar a transformar observaciones de laboratorio en documentación.
\end{itemize}

\subsection*{Reglas de uso de IA durante el taller}

\begin{itemize}
  \item Toda respuesta de IA debe verificarse con una prueba, una gráfica o una lectura de unidades.
  \item La IA no define umbrales normativos; ayuda a razonar alternativas.
  \item Todo cálculo debe tener entrada, unidad, salida y supuesto explícitos.
  \item No se aceptan explicaciones que no se puedan conectar con el flujo real de medición.
  \item Las recomendaciones de IA deben convertirse en entregables concretos: prueba, tabla, diagrama o guía.
\end{itemize}

\section*{Requisitos de software}

\begin{itemize}
  \item Python 3.10 o superior.
  \item Entorno virtual de Python.
  \item Librerías numéricas y científicas para arreglos, señales, datos y gráficas.
  \item Servidor web Python o herramienta equivalente para exponer un servicio de análisis.
  \item Cliente HTTP para enviar tramas de prueba.
  \item Docker opcional para empaquetar servicios.
  \item Git para control de versiones.
  \item Herramienta de IA para apoyo técnico.
  \item SDR opcional. Si no hay hardware disponible, se trabajará con simulación.
\end{itemize}

\section*{Arquitectura conceptual de referencia}

El postproceso se entiende como una cadena de módulos con responsabilidades claras.

\begin{enumerate}
  \item \textbf{Adquisición RF:} antena y SDR reciben señales del entorno.
  \item \textbf{Muestreo IQ:} la señal se representa en componentes I y Q.
  \item \textbf{Buffer temporal:} las muestras se almacenan para que adquisición y procesamiento puedan trabajar a ritmos distintos.
  \item \textbf{Estimación espectral:} las muestras IQ se convierten en una representación de potencia por frecuencia.
  \item \textbf{Correcciones básicas:} se reducen offsets, desbalances y artefactos como el pico DC.
  \item \textbf{Trama de medición:} se empaqueta el espectro con metadatos: rango de frecuencia, timestamp, sensor y configuración.
  \item \textbf{Servicio de análisis:} recibe la trama, estima ruido, detecta emisiones y mide parámetros.
  \item \textbf{Servicio normativo:} compara frecuencia, ancho de banda y potencia contra información autorizada.
  \item \textbf{Respuesta de plataforma:} entrega resultados a backend, interfaz, reportes o almacenamiento.
\end{enumerate}

\section*{Conceptos mínimos de DSP}

El taller usará los siguientes conceptos de forma aplicada:

\begin{itemize}
  \item \textbf{IQ:} representación compleja de una señal RF llevada a banda base.
  \item \textbf{Muestreo:} conversión de una señal continua en valores discretos.
  \item \textbf{FFT:} transformación que permite observar contenido en frecuencia.
  \item \textbf{PSD:} estimación de potencia por frecuencia.
  \item \textbf{RBW:} resolución de frecuencia buscada en el espectro.
  \item \textbf{Ventana:} forma de preparar un bloque antes de FFT para controlar fuga espectral.
  \item \textbf{Promediado espectral:} técnica para reducir variabilidad de la estimación.
  \item \textbf{Banco de filtros:} alternativa con mejor aislamiento entre bins y mayor costo.
  \item \textbf{Piso de ruido:} nivel base del espectro cuando no hay una emisión dominante.
  \item \textbf{Umbral dinámico:} criterio de detección que puede variar según la región del espectro.
  \item \textbf{Ancho de banda ocupado:} rango de frecuencias que concentra la energía relevante de una emisión.
  \item \textbf{SNR:} relación entre señal y ruido.
  \item \textbf{DC:} componente en frecuencia cero, que puede ser artefacto o señal real centrada.
  \item \textbf{PPM:} error relativo del oscilador que desplaza la frecuencia real frente a la nominal.
  \item \textbf{dB y dBm:} escalas logarítmicas para diferencias relativas y niveles de potencia.
\end{itemize}

\section*{Conceptos mínimos de software}

El taller también enseña conceptos de software necesarios para operar y extender el postproceso:

\begin{itemize}
  \item \textbf{Contrato de datos:} estructura mínima que conecta adquisición, análisis y plataforma.
  \item \textbf{API HTTP:} mecanismo para enviar una medición y recibir resultados.
  \item \textbf{Procesamiento por lotes:} análisis de varias tramas con orden y trazabilidad.
  \item \textbf{Configuración de entorno:} rutas, credenciales, parámetros y límites operativos.
  \item \textbf{Validación de entrada:} revisión de campos obligatorios, tamaños, rangos y unidades.
  \item \textbf{Manejo de errores:} respuestas claras cuando faltan datos, hay unidades incorrectas o no hay licencia aplicable.
  \item \textbf{Pruebas sintéticas:} señales diseñadas para saber qué debería detectar el sistema.
  \item \textbf{Observabilidad:} registros, métricas, tiempos de ejecución y muestras de entrada/salida.
  \item \textbf{Empaquetamiento:} ejecución reproducible del servicio en ambiente local o contenedor.
\end{itemize}

\section*{Duración sugerida}

Se propone una duración de 4 semanas con 3 sesiones por semana, cada una de 2 horas. Esto da 24 horas sincrónicas. Se recomienda complementar con 2 a 4 horas semanales de práctica autónoma.

\begin{itemize}
  \item \textbf{Semanas:} 4.
  \item \textbf{Sesiones por semana:} 3.
  \item \textbf{Duración por sesión:} 2 horas.
  \item \textbf{Total sincrónico:} 24 horas.
  \item \textbf{Modalidad:} presencial, remota o mixta.
  \item \textbf{Formato:} explicación breve, laboratorio, discusión técnica y cierre con entregable.
\end{itemize}

\section*{Cronograma detallado}

\begin{longtable}{|p{0.12\textwidth}|p{0.24\textwidth}|p{0.38\textwidth}|p{0.17\textwidth}|}
\hline
\textbf{Semana} & \textbf{Sesión} & \textbf{Contenido y actividad} & \textbf{Entregable} \\
\hline
\endhead

1 & \sesion{1}: visión completa y vocabulario mínimo &
Presentación del flujo señal RF, muestras IQ, espectro, trama de medición, detección, cumplimiento y respuesta de plataforma. Glosario de IQ, FFT, PSD, RBW, ruido, umbral, ancho de banda, DC, PPM, dB y dBm. &
Glosario DSP/software y diagrama general del flujo. \\
\hline

1 & \sesion{2}: señales sintéticas y espectro visual &
Laboratorio para crear señales simples, ruido, tonos cercanos, señal débil y una región ancha. Se calcula y grafica una PSD para construir intuición visual antes de hablar de detección automática. &
Notebook o script con señales sintéticas y gráficas PSD. \\
\hline

1 & \sesion{3}: contrato de medición espectral &
Diseño conceptual de una trama de medición: vector espectral, frecuencia inicial, frecuencia final, timestamp, identificador del sensor, configuración y metadatos. Validación de tamaños, rangos y unidades. &
Contrato de datos de medición y ejemplos JSON conceptuales. \\
\hline

2 & \sesion{4}: adquisición IQ y buffers &
Explicación del flujo del nodo embebido: adquisición continua, muestras IQ, almacenamiento temporal, productor/consumidor y riesgo de pérdida cuando el procesamiento no alcanza el ritmo de captura. &
Diagrama de adquisición con buffers y puntos de riesgo. \\
\hline

2 & \sesion{5}: estimación PSD, resolución y costo &
Relación entre tasa de muestreo, resolución de frecuencia, tamaño de bloque, ventana, solape, promediado y costo computacional. Comparación conceptual entre estimación Welch y banco de filtros. &
Tabla de decisiones PSD: resolución, costo, estabilidad y uso recomendado. \\
\hline

2 & \sesion{6}: correcciones RF y artefactos &
Conceptos de offset DC, desbalance IQ, filtrado de canal, error PPM, saturación, ganancia y pico central. Discusión de cuándo corregir y cuándo conservar datos crudos para inspección. &
Checklist de calidad antes de entregar una trama espectral. \\
\hline

3 & \sesion{7}: servicio de análisis y API &
Diseño de un servicio de análisis: solicitud HTTP, validación de entrada, respuesta JSON, procesamiento por lotes, parámetros configurables y registros de diagnóstico. Pruebas desde cliente HTTP. &
Colección de solicitudes de prueba y formato de respuesta esperado. \\
\hline

3 & \sesion{8}: detección por ruido y umbral dinámico &
Laboratorio con PSD sintética para estimar piso de ruido, construir umbral global, adaptar umbral por regiones, detectar segmentos y separar emisiones cercanas por valles. &
Caso de detección con gráfica de piso de ruido, umbral y regiones. \\
\hline

3 & \sesion{9}: modos operativos de análisis &
Diseño conceptual de tres usos: detectar todas las emisiones, validar picos solicitados y generar evaluación normativa. Se prueban picos presentes, ausentes, débiles y cercanos al margen de frecuencia. &
Matriz de pruebas para detección completa y validación de picos. \\
\hline

4 & \sesion{10}: cumplimiento normativo &
Modelo de cruce con licencias: frecuencia nominal, tolerancia de frecuencia, ancho de banda autorizado, tolerancia de ancho de banda, potencia autorizada, municipio o código territorial y decisión de cumplimiento. &
CSV pequeño de práctica y primer reporte de cumplimiento. \\
\hline

4 & \sesion{11}: potencia, RNI y casos límite &
Diferencia entre nivel pico, potencia integrada, potencia relativa e intensidad de campo estimada. Pruebas con señal débil, licencia ausente, unidad incorrecta, múltiples territorios y datos incompletos. &
Suite manual de casos límite y guía de interpretación de potencia/RNI. \\
\hline

4 & \sesion{12}: integración final y handover &
Prueba de punta a punta con datos simulados o reales: adquisición o simulación, trama espectral, análisis, cumplimiento y respuesta para plataforma. Preparación de guía de operación, diagnóstico y backlog. &
Demo final, runbook y backlog técnico. \\
\hline
\end{longtable}

\section*{Plan por semanas}

\subsection*{Semana 1: intuición DSP y contrato de datos}

\textbf{Objetivo:} construir la intuición mínima para leer un espectro y entender qué debe recibir un servicio de postprocesamiento.

\textbf{Temas principales:}

\begin{itemize}
  \item Flujo completo desde RF hasta resultado normativo.
  \item Diferencia entre señal física, muestras, espectro, trama y respuesta.
  \item Unidades mínimas: Hz, kHz, MHz, dB, dBm y dB sobre ruido.
  \item Señales sintéticas: tono, ruido, dos tonos, emisión débil y banda ancha.
  \item Vector espectral, eje de frecuencia y resolución por bin.
  \item Campos conceptuales de una trama de medición.
  \item Uso de IA para generar ejemplos simples y revisar unidades.
\end{itemize}

\textbf{Laboratorios:}

\begin{itemize}
  \item Construir glosario DSP/software.
  \item Graficar señales simples y sus espectros.
  \item Diseñar una trama JSON conceptual.
  \item Validar tamaños, rangos y unidades de la trama.
  \item Documentar supuestos de frecuencia y potencia.
\end{itemize}

\producto glosario, señales sintéticas, gráficas PSD y contrato conceptual de medición.

\subsection*{Semana 2: nodo de adquisición y generación de espectro}

\textbf{Objetivo:} entender cómo el nodo embebido produce una medición espectral y qué errores o artefactos pueden afectar el postproceso.

\textbf{Temas principales:}

\begin{itemize}
  \item Muestreo IQ y representación compleja.
  \item Flujo productor/consumidor entre adquisición y procesamiento.
  \item Diferencia entre procesamiento de PSD y procesamiento de audio.
  \item Resolución espectral frente a costo computacional.
  \item Ventanas, solape y promediado.
  \item Corrección de DC y desbalance IQ.
  \item Filtrado de canal.
  \item Error PPM y frecuencia nominal frente a frecuencia corregida.
  \item Saturación, ganancia y productos espurios.
  \item Conservación de datos crudos cuando una corrección pueda ocultar contenido real.
\end{itemize}

\textbf{Laboratorios:}

\begin{itemize}
  \item Dibujar el camino de una captura IQ hasta una trama espectral.
  \item Comparar dos resoluciones espectrales con datos sintéticos.
  \item Simular un pico central y discutir si debe corregirse.
  \item Simular saturación o ganancia excesiva en una PSD.
  \item Preparar checklist de calidad de medición.
\end{itemize}

\producto diagrama del nodo de adquisición, tabla de parámetros DSP y checklist de calidad.

\subsection*{Semana 3: servicio de análisis y detección}

\textbf{Objetivo:} diseñar y probar el comportamiento conceptual del servicio de análisis desde una trama espectral hasta una lista de emisiones.

\textbf{Temas principales:}

\begin{itemize}
  \item Servicio de análisis como módulo independiente.
  \item Contrato HTTP: solicitud, respuesta, errores y procesamiento por lotes.
  \item Validación de entrada: vector, frecuencias, timestamp, sensor y configuración.
  \item Corrección opcional de ganancia.
  \item Suavizado de la PSD.
  \item Estimación de piso de ruido global.
  \item Piso de ruido local por regiones.
  \item Umbral dinámico.
  \item Detección de regiones ocupadas.
  \item Separación de emisiones cercanas por valles internos.
  \item Interpretación de resultados de detección.
\end{itemize}

\textbf{Laboratorios:}

\begin{itemize}
  \item Levantar un servicio de análisis o simulador de servicio.
  \item Enviar una trama con una emisión clara.
  \item Enviar una trama sólo con ruido.
  \item Enviar una trama con dos emisiones cercanas.
  \item Cambiar el margen sobre ruido y observar el efecto.
  \item Documentar salidas esperadas y casos ambiguos.
\end{itemize}

\producto colección de solicitudes, gráficas de detección y matriz de comportamiento del detector.

\subsection*{Semana 4: cumplimiento, RNI, pruebas y entrega}

\textbf{Objetivo:} cerrar el flujo normativo, entender las decisiones de cumplimiento y dejar una guía de operación mantenible.

\textbf{Temas principales:}

\begin{itemize}
  \item Base de licencias: frecuencia, ancho de banda, potencia, territorio y estado.
  \item Normalización de unidades.
  \item Coincidencia por frecuencia central.
  \item Tolerancia de frecuencia.
  \item Tolerancia de ancho de banda.
  \item Comparación de potencia.
  \item Decisión de licencia encontrada o no encontrada.
  \item Evaluación por uno o varios territorios.
  \item Filtro de candidatos muy débiles antes de cumplimiento.
  \item Estimación RNI e interpretación de ocupación.
  \item Casos límite y errores esperados.
  \item Empaquetamiento y operación del servicio.
\end{itemize}

\textbf{Laboratorios:}

\begin{itemize}
  \item Crear una base de licencias pequeña para pruebas.
  \item Ejecutar cumplimiento con una emisión que coincide.
  \item Probar licencia ausente.
  \item Probar ancho de banda excedido.
  \item Probar potencia excedida.
  \item Probar múltiples territorios.
  \item Revisar campos de RNI y explicar su alcance.
  \item Preparar demo final y guía de operación.
\end{itemize}

\producto reporte de cumplimiento validado, suite de casos manuales, runbook y backlog.

\section*{Mapa conceptual de módulos}

\begin{longtable}{|p{0.30\textwidth}|p{0.60\textwidth}|}
\hline
\textbf{Módulo conceptual} & \textbf{Qué debe entender el participante} \\
\hline
\endhead

Adquisición RF & Cómo una antena y un SDR generan muestras IQ y qué límites físicos afectan la medición. \\
\hline
Buffer temporal & Por qué adquisición y procesamiento no siempre van al mismo ritmo, y cómo se evitan bloqueos o pérdidas. \\
\hline
Estimación espectral & Cómo se transforma una señal temporal en un espectro de potencia por frecuencia. \\
\hline
Correcciones RF & Qué problemas corrigen DC, balance IQ, filtrado y calibración de frecuencia. \\
\hline
Contrato de medición & Qué información mínima necesita un servicio de análisis para interpretar un espectro. \\
\hline
Servicio de análisis & Cómo validar entrada, estimar ruido, detectar regiones, medir emisiones y responder. \\
\hline
Servicio normativo & Cómo comparar mediciones contra frecuencias, anchos de banda y potencias autorizadas. \\
\hline
Observabilidad & Qué registros, gráficas y casos de prueba ayudan a diagnosticar errores. \\
\hline
\end{longtable}

\section*{Prácticas guiadas sugeridas}

\subsection*{Práctica 1: construir una PSD mínima}

\begin{itemize}
  \item Generar un eje de frecuencia entre dos límites conocidos.
  \item Crear un piso de ruido plano.
  \item Añadir una emisión sintética en una frecuencia central conocida.
  \item Empaquetar la medición como una trama JSON conceptual.
  \item Enviar la trama al servicio de análisis o a un simulador equivalente.
\end{itemize}

\subsection*{Práctica 2: explorar umbral dinámico}

\begin{itemize}
  \item Probar varios márgenes sobre el piso de ruido.
  \item Graficar piso de ruido, umbral y PSD.
  \item Observar cuándo una emisión aparece o desaparece.
  \item Explicar el resultado con unidades y no sólo con intuición visual.
\end{itemize}

\subsection*{Práctica 3: validar picos solicitados}

\begin{itemize}
  \item Solicitar una frecuencia donde sí hay emisión.
  \item Solicitar una frecuencia cercana pero fuera de tolerancia.
  \item Solicitar una frecuencia donde sólo hay ruido.
  \item Solicitar una emisión débil bajo el umbral.
  \item Documentar por qué cada caso se acepta o se rechaza.
\end{itemize}

\subsection*{Práctica 4: cumplimiento con base pequeña}

\begin{itemize}
  \item Crear una tabla pequeña con territorio, frecuencia nominal, ancho de banda autorizado y potencia autorizada.
  \item Probar una emisión que coincide con la autorización.
  \item Cambiar la frecuencia para forzar no coincidencia.
  \item Cambiar el ancho de banda para forzar incumplimiento.
  \item Cambiar la potencia para forzar incumplimiento.
\end{itemize}

\subsection*{Práctica 5: caso de borde RF}

\begin{itemize}
  \item Simular una PSD con energía en el centro exacto.
  \item Discutir si parece artefacto DC o señal real.
  \item Comparar resultado con y sin reparación conceptual del centro.
  \item Decidir si conviene corregir, conservar dato crudo o desplazar sintonía.
\end{itemize}

\section*{Entregables finales}

Al finalizar las 4 semanas se espera entregar:

\begin{itemize}
  \item Glosario DSP/software operativo.
  \item Diagrama del flujo de adquisición, análisis y cumplimiento.
  \item Scripts o notebooks de señales sintéticas.
  \item Contrato conceptual de trama espectral.
  \item Colección de solicitudes de prueba.
  \item Matriz de pruebas para detección completa, validación de picos y cumplimiento.
  \item Base de licencias pequeña para pruebas controladas.
  \item Guía de interpretación de resultados.
  \item Checklist de calidad de medición espectral.
  \item Runbook de operación local o contenedorizada.
  \item Backlog de mejoras técnicas.
\end{itemize}

\section*{Criterios de éxito}

El taller será exitoso si los participantes pueden:

\begin{itemize}
  \item Explicar qué parte ocurre en el nodo de adquisición y qué parte en el servicio de análisis.
  \item Construir una trama espectral válida sin copiar un ejemplo ciegamente.
  \item Enviar una medición a un servicio de análisis y leer la respuesta.
  \item Cambiar un umbral y explicar el efecto.
  \item Interpretar una detección completa de emisiones.
  \item Explicar por qué una frecuencia solicitada se clasifica como emisión o ruido.
  \item Preparar una base de licencias de prueba con unidades correctas.
  \item Explicar cumplimiento de frecuencia, ancho de banda y potencia.
  \item Reconocer riesgos de DC, PPM, ganancia, saturación y potencia relativa.
  \item Usar IA para acelerar análisis sin dejar de validar con pruebas.
\end{itemize}

\section*{Riesgos y mitigaciones}

\begin{itemize}
  \item \textbf{DSP demasiado abstracto:} usar señales sintéticas y gráficas antes de fórmulas.
  \item \textbf{Hardware no disponible:} usar simulación como ruta principal y hardware como extensión.
  \item \textbf{Confusión de unidades:} mantener tabla visible de Hz, MHz, kHz, dB, dBm y dBm/Hz.
  \item \textbf{Potencia mal interpretada:} diferenciar nivel pico, potencia integrada, potencia relativa e intensidad de campo.
  \item \textbf{Dependencia excesiva de implementación existente:} trabajar con módulos conceptuales, contratos y pruebas, no con rutas de código.
  \item \textbf{Cruce normativo confuso:} empezar con una base pequeña y controlada antes de usar bases completas.
  \item \textbf{Exceso de confianza en IA:} exigir gráficas, casos de prueba y explicación de supuestos.
  \item \textbf{Alcance demasiado grande:} cerrar primero el camino mínimo desde trama espectral hasta resultado de cumplimiento.
\end{itemize}

\section*{Agenda sugerida de cada sesión}

Cada sesión de 2 horas puede seguir este formato:

\begin{itemize}
  \item 10 minutos: objetivo y conexión con el flujo completo.
  \item 25 minutos: explicación conceptual con diagrama o gráfica.
  \item 45 minutos: laboratorio práctico.
  \item 20 minutos: revisión de resultados, unidades y supuestos.
  \item 10 minutos: uso de IA para documentar o proponer pruebas.
  \item 10 minutos: cierre y tarea corta.
\end{itemize}

\section*{Material mínimo para el instructor}

\begin{itemize}
  \item Diagrama conceptual del flujo RF a cumplimiento.
  \item Notebook o script para señales sintéticas.
  \item Tres tramas espectrales de ejemplo: ruido, una emisión y varias emisiones.
  \item Base de licencias pequeña de práctica.
  \item Colección de solicitudes HTTP.
  \item Checklist de instalación.
  \item Checklist de demo final.
  \item Prompts guía para IA.
  \item Plantilla de matriz de pruebas.
\end{itemize}

\section*{Backlog técnico sugerido}

Después del taller, el equipo puede continuar con:

\begin{itemize}
  \item Precisar la semántica de potencia usada por cada reporte.
  \item Añadir pruebas automáticas para detección con señales sintéticas.
  \item Crear fixtures de espectros versionados.
  \item Documentar formalmente el contrato entre adquisición y análisis.
  \item Comparar estrategias de detección genérica y detección especializada por banda.
  \item Añadir métricas de rendimiento para análisis por lotes.
  \item Crear tablero de diagnóstico para ruido, umbral y regiones detectadas.
  \item Definir criterios de aceptación para cambios en umbrales y mediciones.
\end{itemize}

\section*{Conclusión}

Este taller de 4 semanas convierte el postprocesamiento en una ruta aprendible: primero se construye intuición DSP mínima, luego se entiende cómo se genera una medición espectral, después se diseña el servicio de análisis y finalmente se cierra el flujo de cumplimiento normativo con pruebas y documentación.

El resultado esperado es un equipo capaz de operar, explicar, probar y extender el postproceso con criterio técnico, sin depender de una comprensión matemática profunda desde el primer día y sin quedar atado a rutas o detalles internos de una implementación específica.

\end{document}
